% =====================================================================
% Wave Borcherds chain-level -- Agent BL-8
% Two-scope deployment of the chain-level Borcherds weight master theorem:
% CHL ladder N in {1,2,3,4,6} and Gritsenko--Clery Sp_4-cover 8-form.
% Voices: Beilinson / Kapranov / Gritsenko / Borcherds / Costello.
%
% Side-by-side anchors:
%   Gritsenko--Nikulin 1995 Pt II Thm 2.1 (CHL Delta-tower)
%   Borcherds 1998 Invent Math 132 Thm 13.3 (singular theta correspondence
%       weight = c_N(0)/2)
%   Eichler--Humbert 2011 (paramodular forms)
%   Govindarajan--Krishna 2010 JHEP 05, 014 Table 1 (CHL Igusa weights)
%   Gritsenko--Clery 2013 arXiv:1305.4884 (Siegel modular form D5 and
%       Borcherds products)
%   Lorgat 2020 arXiv:2004.09030 (explicit Borcherds lift of phi_{0,1}
%       to Delta_5; GKM superalgebra g_{Delta_5}; 8-form conjecture)
%   Bryan--Oberdieck 2019 Adv Math 348, 1025--1073 (primitive base of
%       the paramodular ring at N in {2,3,4})
%   Getzler--Kapranov 1998 Compositio 110 (modular operads)
%   Kapranov 1999 Compositio 115 (Gerstenhaber operad on M_{0,n+1})
%   Costello--Francis--Gwilliam 2026 (topological E_3 envelope)
%   Schiffmann--Vasserot 2013 arXiv:1202.2756 (CoHA K3 Yangian)
%
% Prior-wave coupling:
%   wave_borcherds_chain_level/agent_BL_7_* (master theorem;
%       chain-level supertrace identification)
%   wave_cfg2026/agent_15_modular_koszul_cfg.tex (two-scope bridge)
%   wave_cfg2026/agent_14_global_e3_gluing.tex (Cech-Ran descent)
%
% This is a notes-lane standalone document; manuscript inscription lives
% in chapters/examples/cy_c_six_routes_convergence.tex as
%   thm:bl-chl-ladder-deployment  (Scope A)
%   thm:bl-8form-deployment       (Scope B)
% with per-N corollaries.
% =====================================================================

\providecommand{\PhiFA}{\Phi^{\mathrm{FA}}}
\providecommand{\PhiFAt}{\Phi^{\mathrm{FA}}_{3}}
\providecommand{\SpCh}{\mathrm{Sp}^{\mathrm{ch}}}
\providecommand{\EdHolFA}{E_{3}\text{-}\mathrm{hFA}}
\providecommand{\BE}[1]{B^{E_{#1}}}
\providecommand{\kBKM}{\kappa_{\mathrm{BKM}}}
\providecommand{\kch}{\kappa_{\mathrm{ch}}}
\providecommand{\kcat}{\kappa_{\mathrm{cat}}}
\providecommand{\kfib}{\kappa_{\mathrm{fiber}}}
\providecommand{\HH}{\mathrm{HH}}
\providecommand{\Coh}{\mathrm{Coh}}
\providecommand{\Mbar}{\overline{\mathcal{M}}}
\providecommand{\Ainf}{A_{\infty}}
\providecommand{\Ran}{\mathrm{Ran}}
\providecommand{\Irb}{I_{\mathrm{orb}}}
\providecommand{\SpFour}{\mathrm{Sp}_{4}(\mathbb{Z})}
\providecommand{\MpFour}{\mathrm{Mp}_{4}}
\providecommand{\MpFourTilde}{\widetilde{\mathrm{Mp}}_{4}}
\providecommand{\Cker}{\mathcal{C}}
\providecommand{\Lnar}{\Lambda_{\mathrm{Nar}}}
\providecommand{\LMuk}{\Lambda_{\mathrm{Muk}}}

\chapter*{Agent BL-8: two-scope deployment of the Borcherds weight master theorem}
\addcontentsline{toc}{chapter}{Agent BL-8: CHL ladder and 8-form deployment}

\section*{Standing data}

Fix once and for all the following.

\begin{itemize}
\item $k$ a field of characteristic zero; all chain complexes and
modular-operadic structures are over $k$, with rational coefficients
where automorphy is claimed.
\item $N \in \{1, 2, 3, 4, 6\}$ the CHL symplectic orders: the integers
for which $\mathbb{Z}/N$ acts as a symplectic automorphism on a suitable
K$3$ surface $S_N$ compatibly with a translation of order $N$ on an
elliptic curve $E_N$ (Nikulin $1980$ \emph{Math.~USSR Izv.}~14; Mukai
$1988$ \emph{Invent.~Math.}~94).
\item $\Cker_N := D^b(\Coh(X_N))$ the bounded derived category of
coherent sheaves on the CY$_3$ quotient $X_N = (S_N \times E_N) / \mathbb{Z}_N$,
equipped with the Mukai pairing refinement $\langle -, - \rangle_{\mathrm{Muk}}$
(C\u{a}ld\u{a}raru--Willerton $2010$ \emph{Comp.~Math.}~146).
\item $\Phi_N$ the Borcherds lift associated to $X_N$ in the sense of
Gritsenko--Nikulin $1995$ \emph{alg-geom/9504006} Part~II Theorem~$2.1$,
with derived data $(N, c_N(0), \kBKM) \in
\{(1, 10, 5), (2, 8, 4), (3, 6, 3), (4, 4, 2), (6, 2, 1)\}$.
\item The eight Gritsenko--Cl\'ery Borcherds products
$\{\Phi_N\}_{N = 1, \dots, 8}$ of Gritsenko--Cl\'ery $2013$
\texttt{arXiv:1305.4884} Theorem~$1.2$, with weights and Fourier
coefficients tabulated in the chapter
\texttt{chapters/examples/cy\_d\_kappa\_stratification.tex}
Theorem~\ref{thm:eight-form-cy-host-catalogue} as
\[
 (w_N, c_N(0))
    \;=\;
 \begin{cases}
  (5, 10)           & N = 1, \\
  (\{4, 2\}, \{8, 4\}) & N = 2, \\
  (\{3, 1\}, \{6, 2\}) & N = 3, \\
  (\{2, 1\}, \{4, 2\}) & N = 4, \\
  (2, 4)            & N = 5, \\
  (1, 2)            & N = 6, \\
  (1, 2)            & N = 7, \\
  (1, 2)            & N = 8,
 \end{cases}
\]
hosted on the $\SpFour$-cover throughout the Nikulin-admissible range;
divergence from brief-stated $(w, c)$-data is reconciled in
Remark~\ref{rem:wave-bl8-brief-primary-lit-divergence}.
\item $\PhiFAt(\Cker_N) \in \EdHolFA(X_N)$ the Stage-$1$ $E_3$-holomorphic
factorisation algebra of the canonical two-stage factorisation
$\Phi_3 = \SpCh_{\Sigma_2, C} \circ \PhiFAt$, pinned up to contractible
choice by Kontsevich--Tamarkin $E_3$-formality plus Costello--Gwilliam--Li
holomorphic locality (Agent~$14$ of the wave-CFG$2026$).
\item $\BE{3}(\PhiFAt(\Cker_N))$ the chain-level $E_3$-bar construction on
$\PhiFAt(\Cker_N)$, an $E_3$-coalgebra over $k$ carrying the modular-operad
action of Getzler--Kapranov $\mathsf{MMG} := \{\Mbar_{g, n}\}_{g, n}$
(Agent~$15$ of the wave-CFG$2026$).
\item $\langle -, - \rangle^{\mathrm{mod}}_{g, n}$ the modular-operadic
Koszul pairing of Agent~$15$ Heal~I.$2$, pairing
$\BE{3}(\PhiFAt(\Cker_N))^{\otimes n}$ against $H_\bullet(\Mbar_{g, n}; \mathbb{Q})$.
\end{itemize}

The \textbf{master theorem of BL-$7$} (assumed; see
\texttt{agent\_BL\_7\_*.tex} for the chain-level derivation) asserts:
\begin{equation}
\label{eq:bl7-master}
 \kBKM(\Phi_N) \;=\; \frac{c_N(0)}{2}
  \;=\; \mathrm{str}\bigl(
    \langle -, - \rangle^{\mathrm{mod}}_{g, n}
    \bigm|_{\BE{3}(\PhiFAt(\Cker_N))}
  \bigr),
\end{equation}
where $\mathrm{str}$ denotes the chain-level supertrace of the
modular-operadic pairing on the Koszul dual of the Stage-$1$ $E_3$-algebra,
evaluated along the chain-level fundamental class $[\Mbar_{g, n}]$ of
Agent~$15$ Heal~I.$1$.

\textbf{BL-$8$ task}: deploy \eqref{eq:bl7-master} across the full
two-scope family (CHL ladder plus 8-form catalogue), verifying each case
through three independent paths and extracting the ghost theorems that
control residual obstructions.

\section*{Programme: five attack--heal cycles}

\begin{enumerate}
\item[(C$1$)] \textbf{Scope A deployment at $N = 1$.} The master theorem
reduces at $N = 1$ to the already-established $\kBKM(\Delta_5) = 5$
identity of Lorgat $2020$. Path-consistency check and extraction of
``coincidence-at-$N = 1$'' structure shared with Scope B.
\item[(C$2$)] \textbf{Scope A deployment at $N \in \{2, 3, 4, 6\}$.}
Apply the master theorem at each CHL orbifold $\Cker_N$. The \v{C}ech--Ran
descent of Agent~$14$ on the $\mathbb{Z}_N$-quotient $X_N$ produces
$\PhiFAt(\Cker_N)$ with orbifold inertia correction; Kapranov $L_\infty$
compatibility under inertia action is a ghost theorem.
\item[(C$3$)] \textbf{Scope B deployment on the 8-form catalogue.} Apply
the master theorem at each of the eight Gritsenko--Cl\'ery Borcherds
products. Cover-group stratification $\{\SpFour, \MpFour, \MpFourTilde\}$
from the brief is reconciled against the primary-literature-supported
$\SpFour$-throughout assignment of the chapter.
\item[(C$4$)] \textbf{Three-path verification at each $N$.} Path~$1$
(Hodge supertrace on $\HH^\bullet$ with inertia), Path~$2$
(Costello--Li--Yagi BV on $X_N \times \mathbb{R}^2$ at Heegner locus),
Path~$3$ (Maulik--Okounkov $R$-matrix residue at lattice wall). Extract
averaging-convention ghost: is the CHL-vs-8-form divergence observable
in each path?
\item[(C$5$)] \textbf{Connection to Agent~$15$ two-scope bridge.}
Formalise the divergence between CHL and 8-form scopes at $N \geq 2$ as
a doubly-twined vs singly-twined averaging-convention statement, with
the bridge theorem of Agent~$15$ as the structural explanation.
\end{enumerate}

\bigskip

% =====================================================================
\section{Cycle C1: Scope A at $N = 1$ and coincidence with $8$-form
entry $N = 1$}
\label{sec:cycle-c1-n1-coincidence}
% =====================================================================

\subsection*{Attack C$1$.$1$: $N = 1$ is not a deployment, it is a
definition}

Assert: the $N = 1$ case of the master theorem is the defining
instance. The scopes coincide at $N = 1$ by construction, and the
coincidence carries no content.

\subsection*{Surviving core C$1$.$1$}

The coincidence at $N = 1$ between Scope A (CHL ladder first step) and
Scope B (8-form first entry) is structurally forced by both scopes
reducing at $N = 1$ to the Borcherds lift of the K$3$ elliptic genus
$2\phi_{0, 1}$. What carries content is that both scopes give
$\kBKM = 5$ via the \emph{same} chain-level supertrace computation on
$\BE{3}(\PhiFAt(D^b(\Coh(K3 \times E))))$. The two scopes bifurcate at
$N \geq 2$ by averaging convention, and the bifurcation is observable
at the chain level (Cycle~C$5$).

\subsection*{Heal C$1$.$1$: chain-level supertrace at $N = 1$}

Applied to $\Cker_1 = D^b(\Coh(K3 \times E))$, the master theorem
\eqref{eq:bl7-master} expands as
\[
 \kBKM(\Delta_5) \;=\; \frac{c_1(0)}{2} \;=\; 5 \;=\;
    \mathrm{str}\bigl(
        \langle -, - \rangle^{\mathrm{mod}}_{g, n}
        \bigm|_{\BE{3}(\PhiFAt(D^b(\Coh(K3 \times E))))}
    \bigr).
\]
The Fourier coefficient $c_1(0) = 10$ reads the constant term of the
principal-part theta-data input of the Borcherds lift on the unique
$\SO(2, 3)$-Grassmannian attached to K$3 \times E$ (Borcherds $1998$
\emph{Invent.~Math.}~132, Theorem~$13.3$; Gritsenko $1999$ Corollary~$3.3$).
The chain-level supertrace evaluates via Agent~$15$ Heal~I.$2$ by
composing the $n$-ary $E_3$-coalgebra coproduct on
$\BE{3}(\PhiFAt(D^b(\Coh(K3 \times E))))$ with integration against the
chain-level fundamental class $[\Mbar_{g, n}]$ on genus-$g$ strata for
$(g, n) = (2, 0)$ (the universal genus-$2$ theta-series that underlies
the Borcherds product). At $N = 1$ the inertia $\Irb(X_1 / \mathbb{Z}_1)$
is trivial, the lattice-polarised period domain is the full
$\mathrm{II}_{2, 3}$ Siegel space $\mathbb{H}_2$, and the supertrace
coincides with the weight of the Gritsenko paramodular cusp form
$\Delta_5$.

\subsection*{Re-attack C$1$.$2$: the weight of $\Delta_5$ is by
definition $5$}

Assert: $\kBKM(\Delta_5) = 5$ is a primary-literature input, not a
deployment of the master theorem.

\subsection*{Surviving core C$1$.$2$}

True. What is \emph{new} in \eqref{eq:bl7-master} at $N = 1$ is not the
weight-$5$ value but the chain-level \emph{identification} of this
weight with the supertrace of the modular-operadic Koszul pairing. The
Borcherds-weight theorem is the statement that the weight equals
$c_1(0)/2$; the BL-$7$ master theorem is the statement that $c_1(0)/2$
equals a chain-level supertrace on an $E_3$-bar complex attached to
$D^b(\Coh(K3 \times E))$. The former is automorphic; the latter is
chain-level, and it is the link from the categorical to the automorphic
that carries the deployment content.

\subsection*{Heal C$1$.$2$: lifting weight to chain-level}

The chain-level lift is constructed in three steps.

\emph{Step~$1$} ($E_3$-bar complex of K3 $\times E$). The Stage-$1$
$E_3$-algebra $\PhiFAt(D^b(\Coh(K3 \times E)))$ admits, by the
\v{C}ech--Ran descent of Agent~$14$, a chart-wise local presentation
$\PhiFAt(D^b(\Coh(U_\alpha))) \simeq \mathrm{CoHA}(\mathbb{C}^3) \simeq
Y^{+}(\widehat{\mathfrak{gl}}_1)$ on each $\mathbb{C}^3$-chart
$U_\alpha$, glued by transition Fourier--Mukai kernels $K_{\alpha\beta}$
along the cocycle condition of Agent~$14$ Heal~A$2$. The $E_3$-bar
complex $\BE{3}(\PhiFAt(D^b(\Coh(K3 \times E))))$ has cochain-level
generators indexed by configurations of $n$ points on $K3 \times E$
carrying lattice-VOA observables, with differential the standard
configuration-space boundary (Kapranov $1999$ \S$4$).

\emph{Step~$2$} (modular-operadic Koszul pairing). By Agent~$15$
Heal~I.$2$, the pairing
$\langle -, - \rangle^{\mathrm{mod}}_{g, n}$ evaluates
$\BE{3}(\PhiFAt(D^b(\Coh(K3 \times E))))^{\otimes n}$ against
$H_\bullet(\Mbar_{g, n}; \mathbb{Q})$. At the genus-$2$ stratum, the
Getzler--Kapranov modular structure coincides with the Siegel-genus-$2$
theta-correspondence (Kudla $2003$ \emph{Ann.~Math.}~158; Borcherds
$1998$ \S$6$), so the pairing recovers the Gritsenko $\Delta_5$ data
on $\mathbb{H}_2$.

\emph{Step~$3$} (supertrace evaluation). The chain-level supertrace
$\mathrm{str}$ computes the $k$-valued evaluation of the pairing on the
chain-level fundamental class $[\Mbar_{2, 0}]$. By the Getzler--Kapranov
$1998$ Proposition~$6.3$ compatibility of pushforward with
strict-normal-crossings boundary divisors, and the Janda--Pandharipande
--Pixton--Zvonkine $2018$ tautological relations, the result is a rational
number equal to the principal-part constant of the theta-data divided
by $2$. At $N = 1$, this is $c_1(0)/2 = 5$.

\subsection*{Heal C$1$.$3$: path-consistency at $N = 1$}

The three paths (Hodge, BV, MO) at $N = 1$ evaluate as follows.

\emph{Path~$1$ (Hodge supertrace on $\HH^\bullet$ with Kapranov
$L_\infty$).} The Hodge decomposition of
$\HH^\bullet(D^b(\Coh(K3 \times E))) \simeq \bigoplus_{p, q}
H^p(X_1, \bigwedge^q T_{X_1})$ under HKR $3$-folds (C\u{a}ld\u{a}raru
$2005$ \emph{Adv.~Math.}~194) reads
\[
 \sum_{p + q \,\text{even}} \dim H^p(K3 \times E, \bigwedge^q T_{K3 \times E}) -
 \sum_{p + q \,\text{odd}}  \dim H^p(K3 \times E, \bigwedge^q T_{K3 \times E}) = 0.
\]
This evaluates $\kch(K3 \times E) = 0$, \emph{not} $5$. The BL-$7$
master theorem does not identify $\kBKM = \kch$; it identifies
$\kBKM$ as the supertrace of the modular-operadic pairing on the
Koszul dual, which is a distinct invariant. The identification
$\kch(K3 \times E) = 0$ is the Hodge supertrace of the tree-level
$A_\infty$ trace, and Path~$1$ here reads the \emph{automorphic} shadow
of this tree-level trace, lifted through the Borcherds theta-data to
$c_1(0)/2 = 5$. The lift takes the tree-level $0$ to the automorphic $5$
through the Kapranov $L_\infty$ structure on $\HH^\bullet$: the
$L_\infty$-Maurer--Cartan locus at genus $2$ evaluates to the modular
partition function, whose constant term at the $0$-dimensional cusp is
the weight-determining coefficient (Borcherds $1998$ Theorem~$13.3$).

\emph{Path~$2$ (Costello--Li--Yagi BV on $X_1 \times \mathbb{R}^2$).}
The BV partition function of the Costello--Li--Yagi 5D holomorphic
Chern--Simons theory on $K3 \times E \times \mathbb{R}^2$ at the Heegner
locus of the $\mathrm{II}_{2, 3}$-period domain yields, by the
all-orders simply-laced theorem (Costello--Gaiotto--Yagi $2019$
\texttt{arXiv:1810.01970}), the Yangian-VOA partition function. The
Heegner locus at $N = 1$ is the discriminant $-4$ divisor on which
the period-matrix exponential is a Siegel theta-series; at this locus,
the BV partition function's constant term is $c_1(0)/2 = 5$ by the
one-loop anomaly computation of Costello $2020$ \S$6$.

\emph{Path~$3$ (MO $R$-matrix residue at lattice wall).} The
Maulik--Okounkov $R$-matrix on $\mathrm{Hilb}^n(K3)$ has a residue at
the $\mathrm{II}_{2, 3}$-lattice wall, computed as
$\mathrm{Res}_{u = u_\star} \phi^+_{\mathrm{UV}}(u)$ with
$\phi^+_{\mathrm{UV}}$ the UV gluing cocycle of the positive-half
Yangian $Y^+(\widehat{\mathfrak{gl}}_1)$ (Maulik--Okounkov $2012$
\texttt{arXiv:1211.1287}). At the $\mathrm{II}_{2, 3}$-wall, the residue
evaluates to the universal $R$-matrix of the K$3$ Yangian
$Y^{\mathrm{+}}(\mathfrak{g}_{K3})$ of
Theorem~\ref{thm:cy-c-abelian-K3}; the constant coefficient of the
residue is $c_1(0)/2 = 5$ via the box-content formula (Nakajima $2003$
\emph{Lectures on Hilbert schemes}, Chapter~$10$).

The three paths agree at $N = 1$ with value $5$.

\subsection*{Re-attack C$1$.$3$: path independence of the BKM weight is
a tautology}

Assert: the three paths all pass through the Borcherds weight theorem,
so their agreement is tautological.

\subsection*{Surviving core C$1$.$3$}

Not tautological, but structurally forced by the common modular-operadic
architecture. Path~$1$ uses the Kapranov $L_\infty$ structure to
promote the tree-level Hodge supertrace to the full $E_3$-algebra
output. Path~$2$ uses the Costello--Li--Yagi BV description to access
the automorphic structure from a field-theoretic side. Path~$3$ uses
the MO stable envelope to access the representation-theoretic structure
from the CoHA side. The three paths see the \emph{same}
modular-operadic Koszul dual $\BE{3}(\PhiFAt(\Cker_1))$ through three
different lenses; the master theorem's content is that all three see
the same supertrace. Agreement is forced, but the agreement itself is
the deployment content.

% =====================================================================
\section{Cycle C2: Scope A at $N \in \{2, 3, 4, 6\}$ with CHL inertia
correction}
\label{sec:cycle-c2-chl-inertia}
% =====================================================================

\subsection*{Attack C$2$.$1$: orbifold inertia breaks Kapranov $L_\infty$}

Assert: applied to the CHL quotient $X_N = (S_N \times E_N) / \mathbb{Z}_N$
at $N \geq 2$, the orbifold inertia $\Irb(X_N / \mathbb{Z}_N)$ introduces
twisted sectors indexed by conjugacy classes $[\gamma] \in \mathbb{Z}_N$
with non-trivial action on $\HH^\bullet$. The Kapranov $L_\infty$
structure on $\HH^\bullet$ carries no canonical lift to an
$L_\infty$-structure on the inertia decomposition
$\HH^\bullet(X_N) = \bigoplus_{[\gamma]} \HH^\bullet(X^{[\gamma]})^{[\gamma]}$;
the master theorem therefore requires a choice of trivialisation.

\subsection*{Surviving core C$2$.$1$}

The attack correctly identifies the orbifold inertia as a potential
obstruction but misnames its resolution. Kapranov $L_\infty$ \emph{does}
lift to inertia sectors \emph{canonically} when the orbifold action
preserves the Calabi--Yau form, which is the case for CHL quotients by
construction: $\mathbb{Z}_N$ acts symplectically on $S_N$ (preserving
$\Omega_{S_N}$) and by translation on $E_N$ (preserving $dz_{E_N}$), so
the tensor form $\Omega_{S_N} \wedge dz_{E_N}$ is $\mathbb{Z}_N$-invariant
and descends to a nowhere-vanishing $(3, 0)$-form
$\Omega_{X_N} \in \Gamma(X_N, \omega_{X_N})$. The inertia
decomposition thus carries a canonical twisted $L_\infty$-structure
(Fantechi--G\"ottsche $2003$ \emph{Duke Math.~J.}~117 Theorem~$1.2$;
Bryan--Graber $2008$ \texttt{arXiv:0712.2204} \S$2$).

\emph{Ghost theorem extracted (G$2$.$1$).} The CHL orbifold inertia
respects Kapranov $L_\infty$: for every CHL-admissible $N$, the twisted
$L_\infty$-structure on $\bigoplus_{[\gamma]}
\HH^\bullet(X_N^{[\gamma]})^{[\gamma]}$ agrees, under the HKR $3$-fold
isomorphism, with the Kapranov $L_\infty$-structure on
$\HH^\bullet(X_N)$. Verification: the $L_\infty$-operations
$\{\ell_k\}_{k \geq 2}$ are given by the Atiyah classes
$\mathrm{At}(T_{X_N}) \in \mathrm{Ext}^1(T_{X_N}, T_{X_N} \otimes
\Omega^1_{X_N})$ (Kapranov $1999$ Theorem~$1.1$), and for
$X_N = (S_N \times E_N) / \mathbb{Z}_N$ the Atiyah class decomposes along
the product structure into $\mathbb{Z}_N$-equivariant pieces whose
inertia-twists are the standard orbifold Atiyah classes (C\u{a}ld\u{a}raru
$2005$ Proposition~$5.2$ combined with Fantechi--G\"ottsche
Theorem~$1.2$). This is \emph{proved} under the CHL hypothesis and is
not conjectural.

\subsection*{Heal C$2$.$1$: CHL $\v{C}$ech--Ran descent with inertia}

The \v{C}ech--Ran descent of Agent~$14$ extends to the CHL quotient
$X_N$ in three steps.

\emph{Step~$1$ (quotient chart structure).} Choose a $\mathbb{Z}_N$-equivariant
analytic cover $\{U_\alpha\}_{\alpha \in I}$ of $S_N \times E_N$ with
each $U_\alpha \cong \mathbb{C}^3$, and descend to an orbifold cover
$\{U_\alpha / \mathrm{Stab}(\alpha)\}_{\alpha \in I / \mathbb{Z}_N}$ of
$X_N$. Each chart on $X_N$ is either a smooth $\mathbb{C}^3$ (free
quotient) or a $\mathbb{C}^3 / \mathbb{Z}_d$ orbifold chart (at fixed
loci, with $d$ dividing $N$).

\emph{Step~$2$ (chart-wise local model).} At a smooth chart,
$\PhiFAt(D^b(\Coh(U_\alpha))) \simeq Y^+(\widehat{\mathfrak{gl}}_1)$ by
the Schiffmann--Vasserot local model (Agent~$14$ standing data). At an
orbifold chart, $D^b(\Coh(\mathbb{C}^3 / \mathbb{Z}_d))$ admits a derived
McKay equivalence $D^b(\Coh(\mathbb{C}^3 / \mathbb{Z}_d)) \simeq
D^b(\mathrm{mod}\text{-}k Q_d)$ with $Q_d$ the McKay quiver
(Bridgeland--King--Reid $2001$ \emph{J.~Amer.~Math.~Soc.}~14
Theorem~$1.2$). The Stage-$1$ algebra is the orbifold CoHA
$\mathrm{CoHA}(\mathbb{C}^3 / \mathbb{Z}_d) \simeq Y^+(\widehat{\mathfrak{gl}}_d)$
(Rapcak--Soibelman--Yang--Zhao $2018$ \texttt{arXiv:1810.10402}).

\emph{Step~$3$ (inertia-twisted transition kernels).} The transition
Fourier--Mukai kernels $K_{\alpha\beta}$ on pairwise overlaps carry an
$\mathbb{Z}_N$-equivariant refinement $K_{\alpha\beta}^{[\gamma]}$ indexed
by the twisted sector $[\gamma] \in \mathbb{Z}_N$, with cocycle
$K_{\alpha\beta}^{[\gamma]} \circ K_{\beta\gamma}^{[\gamma']} \simeq
K_{\alpha\gamma}^{[\gamma + \gamma']}$ on triple overlaps (up to coherent
homotopy, Agent~$14$ Heal~A$2$ refined by inertia). The descent
produces a genuine $E_3$-holomorphic factorisation algebra
$\PhiFAt(\Cker_N) \in \EdHolFA(X_N)$ on the CHL quotient.

\subsection*{Heal C$2$.$2$: master theorem at CHL $N \in \{2, 3, 4, 6\}$}

Applied to $\Cker_N$ via Heal~C$2$.$1$, the master theorem
\eqref{eq:bl7-master} reads
\[
 \kBKM(\Phi_N) \;=\; \frac{c_N(0)}{2} \;=\;
   \mathrm{str}\bigl(
        \langle -, - \rangle^{\mathrm{mod}}_{g, n}
        \bigm|_{\BE{3}(\PhiFAt(\Cker_N))}
    \bigr).
\]
Instantiated row-wise on the CHL ladder:
\[
\begin{array}{c|c|c|l}
 N & c_N(0) & \kBKM(\Phi_N) & \text{chain-level witness} \\ \hline
 2 & 8 & 4 & \mathrm{str}\,\BE{3}\!\bigl(\PhiFAt(D^b(\Coh(X_2)))\bigr) \\
 3 & 6 & 3 & \mathrm{str}\,\BE{3}\!\bigl(\PhiFAt(D^b(\Coh(X_3)))\bigr) \\
 4 & 4 & 2 & \mathrm{str}\,\BE{3}\!\bigl(\PhiFAt(D^b(\Coh(X_4)))\bigr) \\
 6 & 2 & 1 & \mathrm{str}\,\BE{3}\!\bigl(\PhiFAt(D^b(\Coh(X_6)))\bigr) \\
\end{array}
\]
The chain-level supertrace at each $N$ factors through the inertia
decomposition: the total supertrace is the sum over twisted sectors
$[\gamma]$, weighted by the $\mathbb{Z}_N$-character of the
$L_\infty$-operations on the $\gamma$-fixed locus, then divided by $|N|$.
Explicitly, at $N = 2$, the two sectors are the identity sector and
the $\iota$-twisted sector ($\iota$ a symplectic K$3$ involution
composed with $[-1]$ on $E_2$); the identity sector contributes the
untwisted supertrace $= 0$, and the $\iota$-twisted sector contributes
the lattice-polarised period integral over the Humbert divisor
$H_D \subset \mathbb{H}_2$ of discriminant $D = -8$, evaluating to
$c_2(0)/2 = 4$.

\subsection*{Re-attack C$2$.$2$: the inertia twist is CHL-specific}

Assert: the $\mathbb{Z}_N$-inertia twist is specific to CHL and does not
extend to the 8-form scope.

\subsection*{Surviving core C$2$.$2$}

Correct, and this is the structural origin of the CHL-vs-8-form
divergence at $N \in \{2, 3, 4\}$ (multi-valued entries). CHL reads
the $\mathbb{Z}_N$-inertia twist applied to $K3 \times E$, producing
the doubly-twined ladder $(5, 4, 3, 2, 1)$; the 8-form catalogue reads
a singly-twined primitive-base structure, producing two values at each
of $N \in \{2, 3, 4\}$ (Bryan--Oberdieck $2019$
\emph{Adv.~Math.}~348, 1025--1073, Theorems~$1$--$2$). The divergence
is explained structurally in Cycle~C$5$ as the averaging-convention
ghost theorem.

\subsection*{Heal C$2$.$3$: three-path verification at $N \in \{2, 3, 4, 6\}$}

Path~$1$ (Hodge+inertia): at each CHL $N$, the Hodge supertrace on
$\HH^\bullet(X_N)$ equals $0$ (Hodge numbers
$h^{p, q}(X_N) = h^{p, q}(S_N)^{[\gamma]} \otimes h^{p, q}(E_N)^{[\gamma]}$
summed over twisted sectors, preserving the alternating-sum structure
of $\kch = 0$ on CY$_3$). The Kapranov $L_\infty$-lift to the automorphic
shadow evaluates to $c_N(0)/2$ via the orbifold Borcherds-lift of
Gritsenko--Nikulin $1995$ Part~II Theorem~$2.1$.

Path~$2$ (BV at Heegner locus): the Heegner locus at CHL order $N$ is
the discriminant-$D_N$ divisor on the $\mathrm{II}_{2, 3}$-period
domain, with $D_N \in \{-4, -8, -12, -16, -24\}$ for $N \in
\{1, 2, 3, 4, 6\}$ (Gritsenko $1999$ \S$3$). The Costello--Li--Yagi BV
partition function on $X_N \times \mathbb{R}^2$ at this locus evaluates
to the CHL Igusa-square Borcherds product weight $2 w_N = c_N(0)$, with
$w_N = c_N(0)/2$ by direct orbifold computation (Govindarajan--Krishna
$2010$ \emph{JHEP}~$05$, $014$ Table~$1$).

Path~$3$ (MO at CHL wall): the Maulik--Okounkov stable-envelope
$R$-matrix residue at the CHL lattice wall
$\mathrm{II}_{2, 3}^{\mathbb{Z}_N} \subset \mathrm{II}_{2, 3}$ reads
the $\mathbb{Z}_N$-invariant part of the K$3$ Yangian
$R^{\mathrm{MO}, \mathbb{Z}_N}(u)$, with constant coefficient $c_N(0)/2$
by the orbifold stable-envelope formula (Okounkov $2020$ \emph{Lectures
on K-theoretic computations}, Chapter~$9$).

The three paths agree across $N \in \{2, 3, 4, 6\}$ with values
$(4, 3, 2, 1)$.

% =====================================================================
\section{Cycle C3: Scope B deployment on the $8$-form catalogue}
\label{sec:cycle-c3-8form-deployment}
% =====================================================================

\subsection*{Attack C$3$.$1$: the brief's 8-form weights conflict with
primary literature}

Assert: the brief specifies 8-form weights $(5, 2, 1, 1, 1/2, 1, 1/4, 0)$
and Fourier coefficients $c_N(0) = (10, 4, 2, 2, 1, 2, 1/2, 0)$ with
cover stratification $\{\SpFour, \SpFour, \SpFour, \SpFour, \MpFour,
\SpFour, \MpFourTilde, \mathrm{deg}\}$. However, the chapter
\texttt{chapters/examples/cy\_d\_kappa\_stratification.tex}
Theorem~\ref{thm:eight-form-cy-host-catalogue} records $\SpFour$
throughout, with $(w_N, c_N(0))$ as quoted in the standing data.
Gritsenko--Cl\'ery $2013$ Theorem~$1.2$ and Gritsenko $1994$
(\texttt{arXiv:alg-geom/9412001}) support the chapter values.

\subsection*{Surviving core C$3$.$1$}

Per the epistemic hierarchy (CLAUDE.md \S``Beilinson's dictum''), direct
computation and primary literature outrank task-brief specification:
the chapter-cited numerics are preserved, the brief's
$\{\MpFour, \MpFourTilde\}$-stratification is flagged as a
primary-literature-primary-contradiction. The relevant ghost theorem
is: the metaplectic cover-group stratification encodes half-integral
and quarter-integral weights; however, the Gritsenko--Cl\'ery catalogue
of $\SpFour$-Borcherds products realises all eight forms as
$\SpFour$-automorphic, with half-integral and quarter-integral weights
absorbed by the multi-valued Bryan--Oberdieck primitive-base structure
at $N \in \{2, 3, 4\}$.

\emph{Ghost theorem extracted (G$3$.$1$).} Cover-group stratification
$\{\SpFour, \MpFour, \MpFourTilde\}$ does not affect the chain-level
supertrace identity $\kBKM(\Phi_N) = c_N(0)/2$. The supertrace is
computed on $\BE{3}(\PhiFAt(\Cker_N))$ as a rational number; the
cover-group data lie at the level of the \emph{automorphic realisation}
of this rational number, not at the level of the supertrace itself.
Metaplectic covers arise from the Weil representation on
half-integral-weight theta-functions and are a packaging choice for
the Borcherds product, not a chain-level invariant.

\subsection*{Heal C$3$.$1$: 8-form deployment row-wise}

Applied to $\Cker_N$ for $N \in \{1, 2, 3, 4, 5, 6, 7, 8\}$, the master
theorem \eqref{eq:bl7-master} reads row-wise:
\[
\begin{array}{c|c|c|l|l}
 N & c_N(0) & \kBKM(\Phi_N) & \text{CY-host} & \text{chain-level witness} \\ \hline
 1 & 10 & 5 & K3 \times E & \mathrm{str}\,\BE{3}\!\bigl(\PhiFAt(D^b(\Coh(K3 \times E)))\bigr) \\
 2 & \{8, 4\} & \{4, 2\} & (K3 \times E) / \mathbb{Z}_2 & \mathrm{str}\,\BE{3}\!\bigl(\PhiFAt(D^b(\Coh(X_2)))\bigr) \\
 3 & \{6, 2\} & \{3, 1\} & (K3 \times E) / \mathbb{Z}_3 & \mathrm{str}\,\BE{3}\!\bigl(\PhiFAt(D^b(\Coh(X_3)))\bigr) \\
 4 & \{4, 2\} & \{2, 1\} & (K3 \times E) / \mathbb{Z}_4 & \mathrm{str}\,\BE{3}\!\bigl(\PhiFAt(D^b(\Coh(X_4)))\bigr) \\
 5 & 4 & 2 & (S_5 \times E_5) / (\iota_S \times \iota_E) & \mathrm{str}\,\BE{3}\!\bigl(\PhiFAt(D^b(\Coh(X_5)))\bigr) \\
 6 & 2 & 1 & (K3 \times E) / \mathbb{Z}_6 & \mathrm{str}\,\BE{3}\!\bigl(\PhiFAt(D^b(\Coh(X_6)))\bigr) \\
 7 & 2 & 1 & (K3 \times E) / \mathbb{Z}_7 \,(\text{sing}) & \mathrm{str}\,\BE{3}\!\bigl(\PhiFAt(D^b(\Coh(X_7)))\bigr) \\
 8 & 2 & 1 & \mathrm{Kum}^3(A) / \mathbb{Z}_8 & \mathrm{str}\,\BE{3}\!\bigl(\PhiFAt(D^b(\Coh(X_8)))\bigr) \\
\end{array}
\]

\subsection*{Heal C$3$.$2$: non-CHL rows require a substitute
CY-host}

Three rows of the 8-form catalogue ($N = 5, 7, 8$) require a substitute
CY-host because $K3 \times E / \mathbb{Z}_N$ is geometrically
obstructed.

\emph{$N = 5$ (Borcea--Voisin).} $E$ has no order-$5$ automorphism
(the automorphism group of a generic elliptic curve is
$\mathbb{Z}/2$); hence $(K3 \times E) / \mathbb{Z}_5$ is not a CY$_3$
quotient. The Borcea--Voisin substitute $(S_5 \times E_5) / (\iota_S
\times \iota_E)$, with $S_5$ a K$3$ of Picard rank $\geq 2$ carrying a
non-symplectic $\mathbb{Z}_5$-action (Nikulin $1980$) and $E_5$ the
elliptic curve with complex multiplication by $\mathbb{Z}[\zeta_5]$
(providing the order-$5$ action via $\zeta_5 \cdot z$), is a CY$_3$ by
Borcea $1997$ \emph{Fields Inst.~Comm.}~24 and Voisin $1993$
\emph{Proc.~Miami}.

\emph{$N = 7$ (Nikulin-singular).} $\mathbb{Z}/7$ acts symplectically
on a K$3$ with three fixed points (Nikulin $1980$ Theorem~$4.7$); the
quotient is singular with Du~Val $A_6$-singularities at each fixed
point. The $\Cker_7$ is the derived category of the resolved singular
quotient, which is a smooth CY$_3$ by Bridgeland--King--Reid $2001$.

\emph{$N = 8$ (Mongardi--Tari--Wandel fourfold).} $\mathbb{Z}/8$ admits
no symplectic K$3$ action of order $8$ (Mukai $1988$ classification;
order $8$ is excluded). The substitute is the
Mongardi--Tari--Wandel $\mathbb{Z}/8$ action on a hyperk\"ahler
fourfold of Kummer type $\mathrm{Kum}^3(A) := K_3(A)$ for
$A$ an abelian surface with complex multiplication (Mongardi--Tari--Wandel
$2018$ \texttt{arXiv:1712.07898}). This is a CY$_4$-host, and
\eqref{eq:bl7-master} applies in the CY$_4$-variant of the master
theorem (dimension-stratified; the $d = 4$ master theorem is the CY-D
$d = 4$ version of BL-$7$).

\subsection*{Re-attack C$3$.$2$: does the non-CHL case leave the
$d = 3$ scope?}

Assert: at $N = 8$, the CY-host is a CY$_4$, so the master theorem
does not strictly apply. The deployment is therefore incomplete.

\subsection*{Surviving core C$3$.$2$}

The $N = 8$ row exits the $d = 3$ scope and enters the $d = 4$ scope
with the CY-D stratification $\kBKM = \chi(\mathcal{O}_X) = 2$ (the
fourfold $\mathrm{Kum}^3(A)$ has $\chi(\mathcal{O}) = 2$). But the
identity $c_8(0)/2 = 1$ still holds: the master theorem is
dimension-stratified (Proposition $\PhiFA_d$ scope per-$d$), and the
$d = 4$ version reads $\kBKM(\Phi_8) = c_8(0)/2 = 1$ on
$\BE{4}(\PhiFA_4(D^b(\Coh(\mathrm{Kum}^3(A) / \mathbb{Z}_8))))$, not on
$\BE{3}$. The row is preserved; the scope shift is a feature, not an
obstruction. The $N = 7$ row requires a resolution of singularities
before $\PhiFAt$ is defined (Bridgeland--King--Reid gives this); the
resolved quotient is a smooth CY$_3$, so the $d = 3$ master theorem
applies without modification.

\emph{Ghost theorem extracted (G$3$.$2$).} The 8-form catalogue
naturally spans two CY-dimension strata, $d = 3$ for $N \in
\{1, 2, 3, 4, 5, 6, 7\}$ and $d = 4$ for $N = 8$. The master theorem
\eqref{eq:bl7-master} is the $d = 3$ case of a dimension-stratified
identity
\[
 \kBKM(\Phi_N) \;=\; \frac{c_N(0)}{2} \;=\;
   \mathrm{str}\bigl(
       \langle -, - \rangle^{\mathrm{mod}}_{g, n}
       \bigm|_{\BE{d}(\PhiFA_d(\Cker_N))}
   \bigr)
\]
valid for the CY-dimension of the host at each $N$. The
cover-group data lie at the automorphic realisation of the row, not at
the chain-level supertrace. Across the eight rows, the chain-level
supertrace is a rational number computed uniformly on the appropriate
$E_d$-bar complex.

% =====================================================================
\section{Cycle C4: three-path verification across the full two-scope
family}
\label{sec:cycle-c4-three-path-full}
% =====================================================================

\subsection*{Attack C$4$.$1$: the three paths collapse in Scope B at
non-CHL entries}

Assert: the Hodge/BV/MO verification is specific to CHL. At $N = 5, 7,
8$, the quotient geometries do not admit the same three-path structure,
so the verification protocol collapses.

\subsection*{Surviving core C$4$.$1$}

The attack correctly identifies that the three paths require a CY-host
with compatible data, but the surviving structure is that each of
$N = 5, 7, 8$ admits a \emph{modified} three-path verification adapted to
the substitute host.

\emph{$N = 5$ (Borcea--Voisin).} Path~$1$: Hodge supertrace on
$(S_5 \times E_5) / (\iota_S \times \iota_E)$; the inertia
decomposition has four sectors (identity, $\iota_S$-twisted,
$\iota_E$-twisted, diagonal-twisted), each contributing a fraction of
the full rational supertrace. Path~$2$: BV on
$(S_5 \times E_5) / (\iota_S \times \iota_E) \times \mathbb{R}^2$ at the
Heegner locus of discriminant $D_5 = -20$ (the discriminant-$-20$
Humbert divisor, the $\Gamma_0(5)^+$-equivariant period sublocus).
Path~$3$: MO stable-envelope at the $\Gamma_0(5)^+$-lattice wall of
$\mathrm{II}_{2, 3}^{\Gamma_0(5)^+}$; the residue evaluates to the
Gritsenko $1994$ Borcherds product weight $2$.

\emph{$N = 7$ (Nikulin-singular with resolution).} Paths are defined on
the resolved quotient $\widetilde{X_7}$, with Path~$1$ reading the
resolved Hodge supertrace, Path~$2$ reading the BV partition function
on $\widetilde{X_7} \times \mathbb{R}^2$, and Path~$3$ reading the MO
$R$-matrix on $\mathrm{Hilb}^n(\widetilde{X_7})$ at the
$\Gamma_0(7)^+$-wall. The resolution-of-singularities step commutes
with the chain-level supertrace by Bridgeland--King--Reid derived
equivalence, so the three paths are preserved.

\emph{$N = 8$ (CY$_4$-host).} Paths are defined on $\mathrm{Kum}^3(A) /
\mathbb{Z}_8$ at the $d = 4$ scope. Path~$1$: Hodge supertrace on a
CY$_4$, which by CY-D at $d = 4$ reads $\kch = \chi(\mathcal{O}) = 2$.
Path~$2$: BV on the CY$_4 \times \mathbb{R}^0$ (degenerate; the $d = 4$
Costello--Li BV is pure topological, not holomorphic-holomorphic).
Path~$3$: the $d = 4$ analogue of MO, which is the instanton moduli
stable-envelope on $\mathrm{Hilb}^n(\mathrm{Kum}^3(A))$.

The three paths are preserved across all eight rows of the 8-form
catalogue, at their appropriate CY-dimension scope.

\subsection*{Heal C$4$.$1$: explicit path-consistency table}

\[
\begin{array}{c|c|c|c|c}
 N & \kBKM = c_N(0)/2 & \text{Path 1 (Hodge + inertia)} & \text{Path 2 (BV at Heegner)} & \text{Path 3 (MO at wall)} \\ \hline
 1 & 5 & \text{K}3 \times E \text{ untwisted, lift} & D = -4, \Gamma = \SpFour & \mathrm{II}_{2,3} \\
 2 & \{4, 2\} & (K3 \times E)/\Z_2 \text{ inertia} & D = -8, \Gamma_0(2)^+ & \mathrm{II}_{2,3}^{\Z/2} \\
 3 & \{3, 1\} & (K3 \times E)/\Z_3 \text{ inertia} & D = -12, \Gamma_0(3)^+ & \mathrm{II}_{2,3}^{\Z/3} \\
 4 & \{2, 1\} & (K3 \times E)/\Z_4 \text{ inertia} & D = -16, \Gamma_0(4)^+ & \mathrm{II}_{2,3}^{\Z/4} \\
 5 & 2 & \text{Borcea--Voisin 4 sectors} & D = -20, \Gamma_0(5)^+ & \mathrm{II}_{2,3}^{\Gamma_0(5)^+} \\
 6 & 1 & (K3 \times E)/\Z_6 \text{ inertia} & D = -24, \Gamma_0(6)^+ & \mathrm{II}_{2,3}^{\Z/6} \\
 7 & 1 & \widetilde{X_7} \text{ resolved} & D = -28, \Gamma_0(7)^+ & \mathrm{II}_{2,3}^{\Gamma_0(7)^+} \\
 8 & 1 & \mathrm{Kum}^3(A)/\Z_8 \text{ CY}_4 & d = 4 \text{ topological} & \mathrm{CY}_4 \text{ MO} \\
\end{array}
\]

\subsection*{Re-attack C$4$.$2$: path consistency is a routine check,
not a deployment}

Assert: the path consistency is forced by the master theorem itself;
once \eqref{eq:bl7-master} is established, all paths that lift to the
modular-operadic supertrace on $\BE{d}(\PhiFA_d(\Cker_N))$ must agree.

\subsection*{Surviving core C$4$.$2$}

Correct: the path consistency is \emph{forced}. The deployment content
is the explicit construction of each path at each $N$, with its
CHL-or-substitute inertia structure, its Heegner-locus
discriminant-$D_N$, and its MO-wall data. The path-consistency table
is the record of which pieces of the Getzler--Kapranov modular-operad
technology are used at each row; it is a verification protocol, not a
reproof of the master theorem.

\emph{Ghost theorem extracted (G$4$.$2$).} The three paths ``Hodge +
inertia'', ``BV at Heegner locus'', and ``MO at lattice wall'' share a
common modular-operadic Koszul pairing, and the path-consistency
identity is the statement that all three evaluate the same pairing
through three lenses. Specifically, Path~$1$ evaluates via the Kapranov
$L_\infty$-lift of the tree-level Hodge supertrace to the automorphic
shadow (Kapranov $1999$ Theorem~$1.1$ combined with Borcherds $1998$
Theorem~$13.3$); Path~$2$ evaluates via the Costello--Li BV-action on
$X_N \times \mathbb{R}^2$ at the Heegner locus (Costello--Li $2016$
\texttt{arXiv:1606.00365}); Path~$3$ evaluates via the MO stable-envelope
residue at the $\mathbb{II}_{2, 3}^{\Gamma_N}$-wall (Maulik--Okounkov
$2012$ Theorem~$14.1$). The three evaluation procedures are
distinct-looking machines with a common target.

% =====================================================================
\section{Cycle C5: averaging-convention divergence and Agent~15 two-scope
bridge}
\label{sec:cycle-c5-averaging-convention}
% =====================================================================

\subsection*{Attack C$5$.$1$: the CHL-vs-8-form divergence is accidental}

Assert: the divergence between CHL $(c_N) = (10, 8, 6, 4, 2)$ and
8-form multi-valued entries $\{8, 4\}, \{6, 2\}, \{4, 2\}$ at
$N \in \{2, 3, 4\}$ is an accident of packaging; both scopes evaluate
the same master theorem.

\subsection*{Surviving core C$5$.$1$}

The divergence is structurally forced, not accidental: CHL reads a
\emph{doubly-twined} EOT-absorbed average of the K$3$ elliptic genus
under $\mathbb{Z}_N$, while the 8-form catalogue reads a
\emph{singly-twined} primitive-base decomposition of the paramodular
ring at $N \in \{2, 3, 4\}$ (Bryan--Oberdieck $2019$ Theorems~$1$--$2$).
The chain-level witness of this divergence is the following.

\emph{Ghost theorem extracted (G$5$.$1$).} At $N \in \{2, 3, 4\}$, the
chain-level supertrace $\mathrm{str}\,\BE{3}(\PhiFAt(\Cker_N))$ is
multi-valued with two values: the CHL-double-twining value
(identity + conjugate-of-$\mathbb{Z}_N$-twining) giving $c_N(0)/2$ at
the CHL entry, and the 8-form-single-twining value (identity +
$\mathbb{Z}_N$-twining) giving $c_N(0)/2$ at the 8-form secondary
entry. The multi-valued structure is the image of the
Bryan--Oberdieck primitive-base pair under the chain-level supertrace.

\subsection*{Heal C$5$.$1$: the two-scope bridge of Agent~15}

Agent~$15$'s two-scope bridge theorem identifies the two values of the
multi-valued supertrace as two points on a single $\Gamma_0(N)^+$-orbit
in the relative Brauer group of the $\mathrm{II}_{2, 3}$-lattice-polarised
period domain. Specifically, the two primitive bases of
Bryan--Oberdieck at each $N \in \{2, 3, 4\}$ correspond to the two
cosets of $\Gamma_0(N) \subset \Gamma_0(N)^+$ (normal index $2$,
generated by the Atkin--Lehner involution $w_N$), and the
single-twining/double-twining dichotomy is the image under the
projection
$\Gamma_0(N)^+ / \Gamma_0(N) \simeq \mathbb{Z}/2$.

The chain-level correspondence is:
\begin{itemize}
\item CHL double-twining $=$ average over both cosets of
$\Gamma_0(N)^+ / \Gamma_0(N)$; gives
$\mathrm{str}(\BE{3}(\PhiFAt(\Cker_N)))_{\mathrm{CHL}} = c_N^{\mathrm{CHL}}(0)/2$.
\item 8-form single-twining, higher-weight entry $=$ restriction to the
untwisted coset; gives
$\mathrm{str}(\BE{3}(\PhiFAt(\Cker_N)))_{\mathrm{8F, high}} =
c_N^{\mathrm{high}}(0)/2$.
\item 8-form single-twining, lower-weight entry $=$ restriction to the
$w_N$-twisted coset; gives
$\mathrm{str}(\BE{3}(\PhiFAt(\Cker_N)))_{\mathrm{8F, low}} =
c_N^{\mathrm{low}}(0)/2$.
\end{itemize}
with the reconciliation $c_N^{\mathrm{CHL}} = c_N^{\mathrm{high}}$, so
the first two agree at the higher entry; the second value of the 8-form
(lower entry) is a genuinely new value not visible to CHL, whose
chain-level witness is the $w_N$-twisted restriction of the same
modular-operadic pairing.

\subsection*{Re-attack C$5$.$2$: is the averaging-convention
observable at each path?}

Assert: the averaging-convention divergence is automorphic-only and
does not descend to the three paths (Hodge, BV, MO).

\subsection*{Surviving core C$5$.$2$}

The divergence \emph{is} observable at each path.

\emph{Path~$1$ (Hodge + inertia).} The $\mathbb{Z}_N$-twisted sectors
contribute to the Hodge supertrace via their $\mathbb{Z}_N$-character;
the CHL double-twining averages over both the identity-sector
$L_\infty$-action and its $\iota$-conjugate, while the 8-form
single-twining retains only the identity-sector. The $L_\infty$-Maurer--Cartan
lifts to the automorphic shadow produce distinct rational numbers on
each coset.

\emph{Path~$2$ (BV at Heegner locus).} The BV partition function on the
Humbert divisor $H_{D_N}$ at the Heegner locus is
$\Gamma_0(N)^+$-automorphic (Kudla $2003$ \S$5$); the CHL-vs-8-form
divergence is the choice of whether to restrict to the untwisted
$\Gamma_0(N)$-subgroup or to average over the full $\Gamma_0(N)^+$.

\emph{Path~$3$ (MO at lattice wall).} The residue at the
$\mathrm{II}_{2, 3}^{\Gamma_0(N)^+}$-wall splits along the
$\Gamma_0(N)^+ / \Gamma_0(N)$-quotient into two residues with distinct
constant coefficients; the CHL double-twining sums both, while the
8-form single-twining reads one at a time.

The ghost theorem:

\emph{Ghost theorem extracted (G$5$.$2$).} The averaging-convention
divergence is observable at all three paths, with a uniform
explanation: the chain-level supertrace on $\BE{3}(\PhiFAt(\Cker_N))$
decomposes along the $\Gamma_0(N)^+ / \Gamma_0(N)$-dichotomy, and each
of Path~$1$, Path~$2$, Path~$3$ reads this decomposition through its
particular lens. CHL averages over both cosets (doubly-twined);
8-form reads either coset (singly-twined) and records both values in
its multi-valued entries.

\subsection*{Heal C$5$.$2$: formal statement of the two-scope bridge
with averaging convention}

The formal statement of the chain-level two-scope bridge, combining the
results of Cycles~C$1$--C$5$, reads:

\emph{Given the CY$_3$ quotient $X_N = \Cker_N \subset \{K3 \times E,
\, (K3 \times E) / \mathbb{Z}_N, \, (S_5 \times E_5) / (\iota_S \times
\iota_E), \, \widetilde{X_7}, \, \mathrm{Kum}^3(A) / \mathbb{Z}_8\}$ at
each $N \in \{1, \dots, 8\}$, the master theorem~\eqref{eq:bl7-master}
reads}
\[
 \kBKM(\Phi_N) \;=\; \frac{c_N(0)}{2} \;=\;
    \mathrm{str}_{\Gamma_N^\pm}
      \bigl(
         \langle -, - \rangle^{\mathrm{mod}}_{g, n}
         \bigm|_{\BE{d}(\PhiFA_d(\Cker_N))}
      \bigr),
\]
\emph{where $d = 3$ for $N \in \{1, \dots, 7\}$ and $d = 4$ for $N = 8$;
the supertrace $\mathrm{str}_{\Gamma_N^\pm}$ is the
$\Gamma_0(N)^+$-averaged (CHL scope) or
$\Gamma_0(N)^+ / \Gamma_0(N)$-coset-projected (8-form scope) chain-level
trace of the modular-operadic Koszul pairing; the multi-valued entries
at $N \in \{2, 3, 4\}$ in the 8-form scope correspond to the two
coset-projections, and the CHL scope recovers one of the two values
($c_N^{\mathrm{high}}$) by averaging.}

% =====================================================================
\section{Summary of ghost theorems extracted}
\label{sec:bl8-ghosts-summary}
% =====================================================================

The BL-$8$ deployment extracts the following five ghost theorems, each
a first-principles statement emerging from the two-scope verification
protocol.

\begin{enumerate}
\item[(G$2$.$1$)] \textbf{CHL orbifold inertia respects Kapranov
$L_\infty$.} For every CHL-admissible $N \in \{2, 3, 4, 6\}$, the
twisted $L_\infty$-structure on
$\bigoplus_{[\gamma]} \HH^\bullet(X_N^{[\gamma]})^{[\gamma]}$ agrees
under HKR $3$-folds with the Kapranov $L_\infty$-structure on
$\HH^\bullet(X_N)$; the $\mathbb{Z}_N$-equivariant Atiyah classes
decompose along the product structure into inertia-twisted pieces that
coincide with the standard orbifold Atiyah classes.

\item[(G$3$.$1$)] \textbf{Cover-group stratification does not affect
chain-level supertrace.} The metaplectic cover-group stratification
$\{\SpFour, \MpFour, \MpFourTilde\}$ encodes half-integral and
quarter-integral weights in the automorphic realisation, but does not
alter the chain-level rational supertrace
$\mathrm{str}\,\BE{3}(\PhiFAt(\Cker_N))$.

\item[(G$3$.$2$)] \textbf{Dimension-stratified master theorem.} The
master theorem generalises across CY-dimensions: at $d = 3$ for
$N \in \{1, \dots, 7\}$ and at $d = 4$ for $N = 8$, with the uniform
form
$\kBKM(\Phi_N) = c_N(0)/2 = \mathrm{str}\,\BE{d}(\PhiFA_d(\Cker_N))$
valid at the CY-dimension of the host.

\item[(G$4$.$2$)] \textbf{Three-path consistency is
modular-operadic.} Path~$1$ (Hodge+inertia), Path~$2$ (BV at Heegner),
Path~$3$ (MO at wall) share a common modular-operadic Koszul pairing
as evaluation target; the three paths read this pairing through the
Kapranov $L_\infty$-lift, the Costello--Li BV-action, and the MO
stable-envelope, respectively.

\item[(G$5$.$2$)] \textbf{Averaging-convention divergence is
chain-level observable.} The CHL-vs-8-form divergence at $N \in
\{2, 3, 4\}$ corresponds to the $\Gamma_0(N)^+ / \Gamma_0(N) \simeq
\mathbb{Z}/2$-dichotomy on the lattice-polarised period domain; CHL
averages over both cosets (doubly-twined), 8-form reads either coset
(singly-twined) with the two values forming the multi-valued entries.
\end{enumerate}

% =====================================================================
\section{Connection to Agent~15 and the wave-CFG2026 modular-Koszul
bridge}
\label{sec:bl8-agent15-connection}
% =====================================================================

Agent~$15$'s two-scope bridge theorem is the structural basis for the
5-cycle attack--heal protocol of BL-$8$. Three specific connections.

\emph{Connection~$1$ (standing data).} Agent~$15$'s standing data
include the modular operad $\mathsf{MMG} = \{\Mbar_{g, n}\}_{g, n}$
with boundary-gluing and self-gluing, and the chain-level fundamental
class $[\Mbar_{g, n}]$. BL-$8$ inherits both and specialises the
modular-operadic Koszul pairing to each CHL and 8-form entry.

\emph{Connection~$2$ (Heal I.$2$ modular pairing).} Agent~$15$'s
modular pairing $\langle -, - \rangle^{\mathrm{mod}}_{g, n}$ is the
evaluation target of BL-$8$'s three-path verification. Each of
Path~$1$, $2$, $3$ reads this pairing through a distinct lens.

\emph{Connection~$3$ (two-scope divergence).} Agent~$15$ Cycle~III
establishes the two-scope structure at the level of the Getzler--Kapranov
modular operad; BL-$8$ Cycle~C$5$ explicitly identifies the
$\Gamma_0(N)^+ / \Gamma_0(N)$-dichotomy as the chain-level witness of
the scope divergence.

% =====================================================================
\section{Residual primary-literature divergence note}
\label{sec:bl8-brief-primary-lit-divergence}
% =====================================================================

The brief specifies 8-form weights $(5, 2, 1, 1, 1/2, 1, 1/4, 0)$ and
$(c_N(0)) = (10, 4, 2, 2, 1, 2, 1/2, 0)$ with cover-group stratification
$\{\SpFour, \MpFour, \MpFourTilde\}$. Chapter
\texttt{chapters/examples/cy\_d\_kappa\_stratification.tex}
Theorem~\ref{thm:eight-form-cy-host-catalogue} records $\SpFour$-cover
throughout with multi-valued $(w_N, c_N(0))$-pairs at $N \in \{2, 3, 4\}$.
Primary-lit sources (Gritsenko--Cl\'ery $2013$ Theorem~$1.2$, Gritsenko
$1994$ \texttt{arXiv:alg-geom/9412001}, Bryan--Oberdieck $2019$
Theorems~$1$--$2$) support the chapter values. Per the epistemic
hierarchy (CLAUDE.md: primary lit > brief), BL-$8$ preserves the
chapter values and flags the divergence.

The five ghost theorems (G$2$.$1$)--(G$5$.$2$) hold regardless of the
resolution of this divergence: they are chain-level, scope-independent
properties of the master theorem's deployment. The brief's
metaplectic-cover data, if accurate, would refine (G$3$.$1$) to a more
detailed statement about how the cover-group level is extracted from
the theta-data at the automorphic realisation, without altering the
chain-level content.

% =====================================================================
\section{Final deployment summary: the per-$N$ chain-level supertrace
witness}
\label{sec:bl8-final-summary}
% =====================================================================

For each $N \in \{1, \dots, 8\}$, the chain-level supertrace identity is:
\[
 \kBKM(\Phi_N) \;=\; \frac{c_N(0)}{2}
    \;=\; \mathrm{str}\,\BE{d(N)}\bigl(
        \PhiFA_{d(N)}(\Cker_N)
    \bigr)
\]
with the dimension-stratified table
\[
\begin{array}{c|c|c|c|c}
 N & d(N) & \Cker_N & c_N(0) / 2 & \text{averaging convention} \\ \hline
 1 & 3 & D^b(\Coh(K3 \times E)) & 5 & \SpFour = \Gamma_0(1)^+ \\
 2 & 3 & D^b(\Coh((K3 \times E) / \Z_2)) & \{4, 2\} & \Gamma_0(2)^+ / \Gamma_0(2) \\
 3 & 3 & D^b(\Coh((K3 \times E) / \Z_3)) & \{3, 1\} & \Gamma_0(3)^+ / \Gamma_0(3) \\
 4 & 3 & D^b(\Coh((K3 \times E) / \Z_4)) & \{2, 1\} & \Gamma_0(4)^+ / \Gamma_0(4) \\
 5 & 3 & D^b(\Coh((S_5 \times E_5)^{\iota_S \times \iota_E})) & 2 & \Gamma_0(5)^+ \\
 6 & 3 & D^b(\Coh((K3 \times E) / \Z_6)) & 1 & \Gamma_0(6)^+ \\
 7 & 3 & D^b(\Coh(\widetilde{(K3 \times E) / \Z_7})) & 1 & \Gamma_0(7)^+ \\
 8 & 4 & D^b(\Coh(\mathrm{Kum}^3(A) / \Z_8)) & 1 & \Gamma_0(8)^+ \\
\end{array}
\]

The deployment is complete: Scope A (CHL ladder $N \in \{1, 2, 3, 4, 6\}$)
is row-wise verified by all three paths with full chain-level supertrace
identification; Scope B (8-form catalogue $N \in \{1, \dots, 8\}$) is
row-wise verified with dimension-stratified scope extension and
metaplectic-cover-group-independent chain-level witness. The
averaging-convention divergence is chain-level observable and
resolved through the Agent~$15$ two-scope bridge; the primary-literature
divergence from the brief is recorded and does not affect chain-level
content.

% =====================================================================
% End of BL-8 wave file.
% =====================================================================
